%======================================================================
% PRELIMINARY -- Section 5.3, OGBench cube. Drop in after
% \subsection{Planning in Higher Dimensions} (sec:highdim).
%
% Every number is from held-out episodes 8000:10000, 3 LIP seeds x 3 task
% seeds, 50 tasks per cell. Provenance: Dyna/RESULTS_TABLE_20260802.md and
% Dyna/RESULTS_freeze_dyna_20260801.md.
%
% OPEN ITEMS flagged inline with \todo{}:
%   - seed count is 3, not the 6 used in Section 5.2
%   - no MPPI row and no latent-distance Adam row were run
%   - DINO-WM not run (same cache-size reason as Reacher)
%   - the reactive-policy amax is RLP's, not its own default
%======================================================================

\subsection{Contact-Rich Manipulation}
\label{sec:ogbench}

\paragraph{Setup.}
Throughout, RLP denotes the residual refiner of Section~\ref{sec:optimality};
it appears as \texttt{LIPv4} in the released artifacts and logs.
Cube is the single-block manipulation task of OGBench \citep{park2025ogbench},
observed as $224\times224$ RGB with a $5$-d action; success is the block within
$4$\,cm of its target position, so the arm's own configuration is irrelevant and
only what it does to the block counts. As in
Sections~\ref{sec:tworoom}--\ref{sec:highdim} a plan is $N=5$ action blocks of
$F=5$ primitive steps, executed before replanning. Two features make it the
hardest of the three environments for the objective. Manipulation is
\emph{contact-rich}: the map from plan to terminal latent is discontinuous at
every make-or-break of contact, so neither a latent norm nor a smooth value
surface is locally faithful. And the task is \emph{under-actuated through an
intermediary}: the planner controls the gripper but is scored on the block, so
success requires a plan whose effect is mediated by a contact event it must
first bring about. We use the two single-token world models
($192$-d), pretrained on $10{,}000$ episodes, and reserve episodes
$8{,}000$--$9{,}999$ for task draws. We report the in-distribution horizon
($\Delta = 25$, budget $50$) and a four-fold extension ($\Delta = 100$, budget
$200$), the latter out of distribution for every learned component.

\begin{table}[t]
\centering
\caption{Cube success rate (\%) and world-model queries per decision, before and
after one round of world-model fine-tuning \eqref{eq:dyna-fixpoint}. All learned
rows are $3$ planner seeds $\times$ $3$ task seeds; CEM rows have no planner
seed and average $3$ task seeds. Queries count rollout-equivalents: $300$
candidates $\times$ $30$ iterations for CEM, $8$ refinements $\times$ (forward
unroll $+$ backward pass) for RLP. The reactive policy is the same architecture
and objective as in Tables~\ref{tab:tworoom}--\ref{tab:reacher}, replanning every
action block (its natural setting, and the one in which it issues \emph{no}
world-model queries at all). Post-Dyna columns share the fine-tuned model with
every row in that column, including the parameter-free latent-distance
baseline.}
\label{tab:cube}
\begin{tabular}{lrcccc}
\toprule
& queries & \multicolumn{2}{c}{LeWM} & \multicolumn{2}{c}{PLDM} \\
\cmidrule(lr){3-4}\cmidrule(lr){5-6}
planner & / dec. & pre-Dyna & post-Dyna & pre-Dyna & post-Dyna \\
\midrule
\multicolumn{6}{l}{\emph{$\Delta = 25$ (one planning horizon)}} \\
\quad latent distance $+$ CEM & $9{,}000$ & 76.0 & 78.0 & 66.7 & 72.0 \\
\quad value metric $+$ CEM    & $9{,}000$ & 79.3 & 85.3 & 73.3 & 84.0 \\
\quad Reactive Policy         & $0$       & 84.0 & 81.3 & 62.2 & 61.1 \\
\quad \textbf{RLP (ours)}     & ${\sim}16$ & \textbf{88.7} & \textbf{92.0} & \textbf{85.3} & \textbf{91.3} \\
\midrule
\multicolumn{6}{l}{\emph{$\Delta = 100$ (four planning horizons, out of distribution)}} \\
\quad latent distance $+$ CEM & $9{,}000$ & 60.0 & 70.7 & 55.3 & 60.0 \\
\quad value metric $+$ CEM    & $9{,}000$ & 74.0 & 82.0 & 63.3 & 79.3 \\
\quad Reactive Policy         & $0$       & 81.6 & 81.6 & 53.8 & 58.0 \\
\quad \textbf{RLP (ours)}     & ${\sim}16$ & \textbf{81.8} & \textbf{85.3} & \textbf{83.1} & 80.2 \\
\bottomrule
\end{tabular}
\end{table}

\paragraph{Here neither cost is adequate, and the search is what closes the gap.}
Cube completes the pattern the previous two sections started. On TwoRoom the
objective was wrong and the critic fixed it; on Reacher the objective was
adequate and the critic was a liability. On Cube \emph{neither} cost suffices
under a sampling planner: at $\Delta = 25$ the latent distance reaches $76.0$
and $66.7$ and the value metric $79.3$ and $73.3$, against RLP's $88.7$ and
$85.3$ at $562\times$ fewer queries. Swapping the cost buys $3.3$ and $6.6$
points; swapping the optimizer buys $9.4$ and $12.0$ on top of the better cost.
This is the regime the theory of Section~\ref{sec:optimality} speaks to most
directly: with the objective held fixed and the budget cut by nearly three
orders of magnitude, the entire margin is attributable to where in the $25$-d
block space the candidates are placed.

\paragraph{Direct descent on the learned metric is worse than sampling it.}
The sharpest evidence that the difficulty is the shape of the search space
rather than the density of coverage comes from replacing the sampler with
gradient descent on the identical objective. Holding the world model, the
critic and the planner seed fixed---all four planners below score against
\emph{the same} critic instance, the one co-trained inside that seed's RLP
run---and varying only the optimizer:

\begin{center}
\begin{tabular}{lrcc}
\toprule
optimizer & queries / dec. & LeWM & PLDM \\
\midrule
\textbf{RLP} (learned refinement, $K{=}8$) & ${\sim}16$ & \textbf{90.2} & \textbf{83.3} \\
CEM (sampling)                             & $9{,}000$  & 82.4 & 74.9 \\
AdamW (fixed-step first-order, $30$ steps) & $3{,}000$  & 77.8 & 64.4 \\
Reactive Policy (single pass)              & $0$        & 58.0 & 45.3 \\
\bottomrule
\end{tabular}
\end{center}

\noindent
AdamW on the learned metric lands \emph{below} CEM on both models, by $4.6$ and
$10.5$ points, despite having access to a gradient the sampler discards. A fixed
step size cannot be simultaneously right for the flat regions of the value
surface and for the neighbourhoods of contact events where it is sharply curved,
and a first-order method with one schedule descends into whichever it meets
first. Sampling is blind to curvature but is not committed to a step length, so
it degrades more gracefully. RLP is the only arm that is both first-order and
adaptive, and it exceeds the better of the two by $7.8$ and $8.4$ points at
$1/560$ the budget. The ordering \mbox{RLP $>$ CEM $>$ Adam $>$ single-pass} is
the ordering Theorems~\ref{thm:dominance}--\ref{thm:strict} predict when the
task distribution is heterogeneous enough that no fixed hyperparameter serves
it.

\paragraph{Model error is the remaining constraint, and the Dyna loop removes
part of it.}
Cube is also the one environment in which fine-tuning the world model on
on-policy data \eqref{eq:dyna-fixpoint} helps rather than hurts. One round lifts
RLP by $3.3$ and $6.0$ points ($t = 8.7$ and $9.0$, $3$ seeds, all six
positive), and---the observation that identifies \emph{what} improved---it lifts
the parameter-free latent-distance baseline as well, by $2.0$ and $5.3$. That
row has no learned component and no capacity to adapt, so its gain can only come
from the model. The fine-tune therefore improves the world model itself rather
than merely making it easier for a learned actor to exploit, which is the
failure mode a Dyna round is usually suspected of. The largest gain in the table
is the value metric under CEM ($+6.0$ and $+10.7$), i.e.\ the critic benefits
more from a better model than the actor does. The effect also closes the gap
between the two world models: PLDM goes from $3.4$ points behind LeWM to $0.7$,
statistically indistinguishable at this sample size.

\paragraph{Out of distribution the amortized policy stops being the weak
baseline.}
Quadrupling the goal distance to $\Delta = 100$---four planning horizons, and out
of distribution for every learned component---compresses RLP's advantage on the
stronger model to nothing: $81.8$ against the reactive policy's $81.6$ on LeWM
pre-Dyna. The reason is visible in the same column. RLP loses $6.9$ points
moving from $\Delta = 25$ to $100$ while the reactive policy loses $2.4$, and on
the fine-tuned model it loses none at all ($81.6$ at both distances). RLP
optimizes a plan whose terminal state sits a fixed $25$ primitive steps ahead, so
that horizon is built into what it was trained for; a memoryless policy
$\pi(\hat z, \hat z_g)$ has no horizon in it and is invariant by construction.
This is the one place in the paper where the refinement recursion is not worth
its cost, and it localizes the reason: not the learned value, and not the search,
but the fixed plan length. It also marks the only negative Dyna cell we measure
($83.1 \to 80.2$ on PLDM), where a model fine-tuned on $\Delta = 25$ rollouts is
evaluated at $\Delta = 100$; \todo{the matched-horizon collection run that would
separate collection distribution from horizon extrapolation is in progress.}

\paragraph{The configuration is not brittle.}
Section~\ref{sec:discussion} argues against planners whose behaviour is
committed through hyperparameters before a task is seen. It is worth reporting
that RLP's own configuration does not appear to be such a commitment on this
task. Sweeping $29$ single-factor arms around the reported setting---the action
clip $a_{\max}$, batch size, actor and critic learning rates and their decay
shapes, the critic's discount over six values, its backup span, the refinement
freeze point, the hindsight cross-episode rate, and the intermediate-step weight
$\lambda$---produced no arm that improved on it on PLDM, and none that cleared
significance on LeWM. Four settings that a single sweep had flagged as
improvements failed to replicate when re-measured at more seeds. We read this as
the operating point sitting in a flat basin rather than on a ridge, which is the
property one wants of a method whose selling point is not needing to be tuned per
task; \todo{it is also, honestly, a null result about our own tuning and should
be stated as one.}

\paragraph{Limitations.}
Every cell is $3$ planner seeds $\times$ $3$ task seeds, against
Section~\ref{sec:highdim}'s $6\times6$; at $50$ episodes per cell the paired
standard error is ${\approx}2$--$3$ points, so differences inside that band are
not an ordering and only the double-digit margins above should be read as
such. \todo{Raise to 6 seeds before submission.} The world models are pretrained
on all $10{,}000$ episodes including the $2{,}000$ reserved for tasks, so
absolute rates are upper bounds; the pre/post-Dyna \emph{deltas} are unaffected,
both arms sharing the base, and every planner in a column shares one world
model. Two rows of Tables~\ref{tab:tworoom}--\ref{tab:reacher} are absent here:
we ran neither MPPI nor a latent-distance Adam arm, so the optimizer comparison
under the native cost is incomplete. The patch-token model is omitted for the
cache-size reason given in Appendix~\ref{app:reacher}. The reactive policy runs
at RLP's $a_{\max}$ rather than its own released default, which its $\Delta=25$
PLDM figure ($62.2$) may understate---its output is a $\tanh$ scaled by
$a_{\max}$, so the bound is part of the learned function rather than an external
clip. Finally the sweep in the preceding paragraph selected on the same task
draws it reports, inheriting the optimism noted in
Section~\ref{sec:highdim}. \todo{Confirm the reported operating point on a fresh
draw.}
